\section{Introduction}
\label{sec:intro}

\subsection{Context}
\label{ssec:context}

The Chinese Remainder Theorem (CRT) identifies, for any squarefree
integer $m = p_1 \cdots p_k$, the residue ring
$\mathbb{Z}/m\mathbb{Z}$ with the product
$\prod_{i=1}^{k} \mathbb{Z}/p_i\mathbb{Z}$ as a ring isomorphism. The
identification is purely algebraic, but it has dynamical content: any
Markov chain on $\mathbb{Z}/m\mathbb{Z}$ whose transition law respects
the coprimality structure of the moduli factorises as a tensor product
of chains on the prime components~\cite{Seneta2006,Walters1982}. The
present paper studies one such chain, the prime-sieve transfer matrix
of Persistence Theory~\cite{PT_MATHEMATICS_M1,PT_Monograph}, and asks
whether the factorisation property survives at the level of mutual
information.

The prime-sieve transfer matrix $T_m$ is defined on the residues of
prime gaps $g_n = p_{n+1} - p_n$ modulo a squarefree~$m$, and it
admits a tensor decomposition $T_m = \bigotimes_{p \mid m} T_p$
inherited from CRT (see Theorem~\ref{thm:crt_tensor}). Its unique
left Perron eigenvector, the \emph{Ruelle measure}
$\pi_m^{\mathrm{Ruelle}}$, factorises as
$\bigotimes_{p \mid m} \pi_p^{\mathrm{Ruelle}}$ and therefore exhibits
exact informational independence between modular channels. The
empirical stationary measure $\pi_m^{\mathrm{emp}}$ of the actual
prime gap sequence, by contrast, exhibits substantial mutual
information between distinct active channels: the multi-modular
predictor of~\cite[Ch.~math\_structures]{PT_Monograph} reports
$\mathrm{MI}(\bmod\,3, \bmod\,5, \bmod\,7) \sim 1$ nat in the
aggregate, and the more refined pairwise values
$\mathrm{MI}(r^{(3)}; r^{(5)}) = 0.138$ bits and
$\mathrm{MI}(r^{(3)}; r^{(7)}) = 0.283$ bits in
ch.~9~\cite[S15.6.259]{PT_Monograph}. These two facts --- exact
independence of the Ruelle measure, substantial mutual information of
the empirical measure --- raise an obvious tension: which of the two
measures governs the physically meaningful sieve dynamics, and what
exactly is the difference between them?

A second motivation comes from the gravitational sector of Persistence
Theory. Above the Hartle--Hawking transition $\mu_c \approx 6.97$,
the spectral data of $T_m$ at the canonical fixed point
$\mu^\ast = 15$ assemble into a Bianchi-I Lorentzian metric
$G(\mu) = \mathrm{diag}(-1, a_3^2, a_5^2, a_7^2)$, where
$a_p(\mu) = \gamp{p}(\mu) / \mu$ is the per-channel scale factor
derived from the holonomy theorem
T6~\cite[Ch.~6,~12]{PT_Monograph}. The emergent space-time of the
theory therefore has the three active primes $\{3, 5, 7\}$ as its
three spatial directions, and the question of inter-channel
correlations becomes a question about a hypothetical spatial coupling
in the emergent geometry. If the empirical correlations were
generated by the metric $G(\mu)$ rather than by the underlying
arithmetic, the theory would predict a $\mu$-dependent mutual
information --- in particular a measurable signature of the
Hartle--Hawking transition in the residual
$\mathcal{I}_{p,q}(\mu)$. The two theorems of this paper rule out
that scenario.

A third motivation is foundational. The companion
paper~\cite{PT_HOLONOMY_NONLOCALITY} argues that the observed
correlations between space-like separated detectors in Bell-type
experiments~\cite{Bell1964,AspectGrangier1982,Hensen2015} are
inherited from the modular structure of the prime sieve, rather than
transmitted through the light-cone of the emergent geometry. That
argument relies on the precise statement, to be proved here, that
$G(\mu)$ does not couple the disjoint modular fibres. The present
paper provides the measure-theoretic and information-theoretic
machinery needed to make that statement rigorous.

\subsection{Main results}
\label{ssec:main_results}

The three central results of this paper are stated informally
here; precise statements appear in
Section~\ref{sec:proofs}.

\begin{description}[itemsep=4pt,leftmargin=*]
  \item[\textbf{Theorem~\ref{thm:empirical_invariance} (empirical form).}]
  For any pair of distinct active primes $p, q \in \{3, 5, 7\}$
  and observables $\Phi_p \in \mathcal{A}_p$, $\Phi_q \in
  \mathcal{A}_q$ on the corresponding modular $\sigma$-algebras,
  the function
  $\mu \mapsto I_{\pi_m^{\mathrm{emp}}}(\Phi_p\,;\,\Phi_q\,|\,
  G(\mu))$ is constant on the emergent Lorentzian domain
  $\mu \in [\mu_c, \infty)$.

  \item[\textbf{Theorem~\ref{thm:ruelle_zero} (idealised form).}]
  Under the Ruelle measure $\pi_m^{\mathrm{Ruelle}}$ of the
  transfer matrix $T_m$, we have
  $I_{\pi_m^{\mathrm{Ruelle}}}(\Phi_p\,;\,\Phi_q) = 0$ exactly.

  \item[\textbf{Corollary~\ref{cor:closed_form} (closed form).}]
  The empirical inter-channel mutual information admits the
  arithmetic closed-form expression
  $\mathcal{I}_{p,q} = (2 \ln 2)^{-1}[\cos^{2k}\theta_3 +
  \sin^{4k}\theta_2]$, with profile depth $k = k_{p,q}$ determined
  by the CRT cascade.
\end{description}

The first theorem is the geometric-invariance statement: the
emergent Bianchi-I metric does not source inter-channel mutual
information, no matter where one is on its Lorentzian branch. The
second theorem is the idealised-vanishing statement: were one to
replace the empirical gap field with a Markov chain whose stationary
measure is the Ruelle one, the inter-channel mutual information
would be zero. The corollary fits the residual against the corpus
data and exhibits the value-by-value match. The four-step proof of
Theorem~\ref{thm:empirical_invariance} consists of the following
ingredients, each of which is conceptually independent:
\begin{enumerate}[label=(\roman*),itemsep=2pt]
  \item the empirical measure of an arithmetically defined trajectory
    cannot depend on a parameter that does not enter the
    construction of the trajectory;
  \item the metric $G(\mu)$ is, by direct functional dependence on
    the spectrum of $T_m$, measurable with respect to the
    shift-invariant $\sigma$-algebra on the path space;
  \item the prime-sieve flow is ergodic, hence the invariant
    $\sigma$-algebra is trivial modulo the empirical measure, hence
    every shift-invariant random variable is almost surely
    constant;
  \item conditioning a joint distribution on an almost surely
    constant random variable does not change its mutual information.
\end{enumerate}
The first ingredient is the central conceptual move of the paper:
it is the formal counterpart of the intuition that the prime gaps
are arithmetic, not geometric. The remaining three are standard
applications of ergodic theory and information theory.

\subsection{Organisation}
\label{ssec:organisation}

Section~\ref{sec:setup} fixes notation and introduces the two
stationary measures $\pi_m^{\mathrm{Ruelle}}$ and
$\pi_m^{\mathrm{emp}}$, the transfer matrix $T_m$ and the emergent
geometry $G(\mu)$.
Section~\ref{sec:crt_factor} establishes the CRT tensor
factorisation of $T_m$ and the corresponding factorisation of
$\pi_m^{\mathrm{Ruelle}}$, and reports the numerical verification
$\|\pi_{15}^{\mathrm{Ruelle}} - \pi_3^{\mathrm{Ruelle}} \otimes
\pi_5^{\mathrm{Ruelle}}\|_2 = 5.88 \times 10^{-53}$ at fifty digits
of precision.
Section~\ref{sec:geometric_fiber} sets up the geometric side: the
holonomy angles $\theta_p$, the anomalous dimensions $\gamp{p}$, and
the family $G(\mu)$ in terms of the spectrum of $T_m$;
Lemma~\ref{lem:G_measurable} expresses $G(\mu)$ as a measurable
function of $\mathrm{Spec}(T_m)$ alone.
Section~\ref{sec:ergodic} derives the ergodicity of the sieve flow
from the spectral gap of $T_m$ and Mertens' theorem, and exhibits
the triviality of the invariant $\sigma$-algebra
(Corollary~\ref{cor:trivial}).
Section~\ref{sec:proofs} assembles the four-step proof of
Theorem~\ref{thm:empirical_invariance} and the algebraic proof of
Theorem~\ref{thm:ruelle_zero}.
Section~\ref{sec:closed_form} derives the closed form for the
empirical residual $\mathcal{I}_{p,q}$ and validates it against the
corpus data; the open question of a first-principles derivation of
the depth parameter $k_{p,q}$ is flagged in
Remark~\ref{rmk:k_interp}.
Section~\ref{sec:applications} discusses the implications, with
particular attention to the foundational reading of the result for
Bell-type non-locality (deferred in detail to the companion
paper~\cite{PT_HOLONOMY_NONLOCALITY}).
Appendix~\ref{app:perron_frobenius} collects the Perron--Frobenius
facts used in the main text;
Appendix~\ref{app:numerics} reports the full numerical output of
the validation script \texttt{scripts/test\_crt\_decoupling.py}.

\subsection{Notation}
\label{ssec:notation}

We collect the symbols used throughout the paper.
\begin{itemize}[itemsep=2pt,leftmargin=*]
  \item $p, q, p_i \in \{2, 3, 5, 7, \ldots\}$ denote rational primes;
    by default we work with the \emph{active} primes
    $p \in \{3, 5, 7\}$ at the canonical fixed point
    $\mustar = 15$.
  \item $g_n = p_{n+1} - p_n$ is the $n$-th prime gap;
    $r_n^{(m)} = g_n \bmod m$ is its mod-$m$ residue.
  \item $\mathcal{S}_p = \{1, \ldots, p-1\}$ is the active subspace
    of $\mathbb{Z}/p\mathbb{Z}$, excluding the class $0$ that is
    forbidden by Theorem~T1.
  \item $T_p$ is the restricted transfer matrix on $\mathcal{S}_p$;
    $T_m = \bigotimes_{p \mid m} T_p$ for $m$ squarefree.
  \item $\pi_p^{\mathrm{Ruelle}}$ is the left Perron eigenvector of
    $T_p$, normalised to a probability;
    $\pi_p^{\mathrm{emp}}$ is the long-time empirical measure of
    $\{r_n^{(p)}\}_n$.
  \item $\sinp{p}(\mu) = \delta_p(2 - \delta_p)$ with
    $\delta_p = (1 - q^p)/p$ is the holonomy
    identity (theorem T6); $\gamp{p}(\mu) = -\partial \ln \sinp{p}/\partial \ln \mu$
    is the corresponding anomalous dimension.
  \item $\qplus = 1 - 2/\mu$ and $q_- = e^{-1/\mu}$ are the two
    bifurcation values of $q$; throughout this paper we use
    $\qplus$ unless explicitly stated.
  \item $a_p(\mu) = \gamp{p}(\mu)/\mu$ is the per-channel scale
    factor;
    $G(\mu) = \mathrm{diag}(-1, a_3^2, a_5^2, a_7^2)$ is the
    Bianchi-I metric on the active primes;
    $\mu_c \approx 6.97$ is the Hartle--Hawking transition,
    $\mustar = 15$ is the canonical fixed point.
  \item $\mathcal{A}_p = \sigma(\{r_n^{(p)} : n \in \mathbb{N}\})$ is
    the $\sigma$-algebra of channel-$p$ observables on the path
    space $\Omega$;
    $\Phi_p$ denotes a bounded $\mathcal{A}_p$-measurable function.
  \item $I_{\nu}(\,\cdot\,;\,\cdot\,)$ denotes mutual information
    under the measure $\nu$, in bits unless stated; $\DKL$ denotes
    Kullback--Leibler divergence;
    $\|\cdot\|_{\mathrm{TV}}$ is total variation.
  \item $\mathcal{I}_{p,q}$ is the empirical inter-channel residual
    of Eq.~\eqref{eq:residual_def}.
\end{itemize}

All numerical values quoted in the paper are computed at
\texttt{mpmath} precision $50$~digits, and are reproducible via the
script \texttt{scripts/test\_crt\_decoupling.py}.
